% Beamer+ demo deck — demo/slides.pdf
%
% Rebuilt from the compiled PDF (the original source was lost). Every slide is
% laid out on a fixed grid so the Beamer+ demo config — videos, 3D models,
% widgets, reveal boxes and ink, all placed by position on the page — still
% lands where it did.
%
%   Coordinates: inside a \begin{page}…\end{page} picture, (x,y) is in pt from
%   the TOP-LEFT corner of the 362.835 × 272.126 pt page, y pointing DOWN.
%
% Build:  pdflatex demo.tex   (twice: the pictures are placed on the page)
%         then copy demo.pdf to demo/slides.pdf.

\documentclass[11pt,aspectratio=43]{beamer}

\usefonttheme[onlymath]{serif}
\usepackage{amsmath,amssymb}
\usepackage{tikz}
\usetikzlibrary{arrows.meta,calc,decorations.pathmorphing}
\usepackage{graphicx}

% ── Colours ────────────────────────────────────────────────────────────────
\definecolor{hl}{RGB}{248,196,69}        % highlighter yellow
\definecolor{cap}{RGB}{127,128,127}      % grey captions
\definecolor{ink}{gray}{0.3}             % icon ink
\definecolor{cardbg}{RGB}{247,247,245}
\definecolor{cardline}{RGB}{217,217,214}
\definecolor{rec}{RGB}{229,72,77}        % record red
\definecolor{deoxy}{RGB}{158,158,255}
\definecolor{oxy}{RGB}{255,158,158}

% ── Beamer chrome: no navigation, a three-part grey footline ───────────────
\setbeamertemplate{navigation symbols}{}
\setbeamertemplate{background canvas}[default]
\setbeamercolor{background canvas}{bg=white}
\setbeamertemplate{footline}{%
  \leavevmode\hbox{%
  \begin{beamercolorbox}[wd=.333333\paperwidth,ht=5.79pt,dp=2.85pt,center]{foot1}%
    \tiny A Guided Tour\end{beamercolorbox}%
  \begin{beamercolorbox}[wd=.333333\paperwidth,ht=5.79pt,dp=2.85pt,center]{foot2}%
    \tiny Beamer$^+$\end{beamercolorbox}%
  \begin{beamercolorbox}[wd=.333333\paperwidth,ht=5.79pt,dp=2.85pt,right]{foot3}%
    \tiny Interactive Presentations\hspace*{2em}\insertframenumber{} / \inserttotalframenumber\hspace*{2ex}\end{beamercolorbox}}}
\setbeamercolor{foot1}{fg=black,bg=black!30}
\setbeamercolor{foot2}{fg=black,bg=black!25}
\setbeamercolor{foot3}{fg=black,bg=black!20}

% ── Placing things on the page ─────────────────────────────────────────────
% (x,y) in pt from the top-left corner, y down.
\newenvironment{page}{%
  \begin{tikzpicture}[remember picture,overlay,
      shift={(current page.north west)},x=1pt,y=-1pt,
      every node/.style={inner sep=0pt,outer sep=0pt}]}%
  {\end{tikzpicture}}
% \txt[node options]{x}{baseline y}{text}
\newcommand{\txt}[4][]{\node[anchor=base west,#1] at (#2,#3) {#4};}
\newcommand{\ctxt}[4][]{\node[anchor=base,#1] at (#2,#3) {#4};}

% The slide's header: the title on a highlighter, a grey line under it and a
% smaller "example usage" line. #1 = baseline of the title.
\newcommand{\head}[4]{%
  \node[anchor=base west,fill=hl,font=\Large\bfseries,
        text height=12.62pt,text depth=3.46pt] at (10.91,#1) {#2};
  \txt[font=\normalsize,text=cap]{10.91}{#1+16.05}{#3}
  \txt[font=\footnotesize,text=cap]{10.91}{#1+29.72}{#4}}

% A keyboard key with an arrow on it, sitting on the text baseline.
\newcommand{\arrowkey}{%
  \tikz[baseline=0pt,x=1pt,y=1pt]{%
    \draw[fill=black!4,draw=black!55,line width=.45pt,rounded corners=1.6pt]
      (0,-1.49) rectangle (11.42,9.11);
    \draw[black!35,line width=.35pt] (1.1,-1.49) -- (10.32,-1.49);
    \draw[black!80,line width=.6pt,-{Stealth[length=2.1pt,width=2.6pt]}] (2.76,3.76) -- (8.46,3.76);}}

% "try …!" hints: a small hand-drawn arrow ending at #3,#4 from #1,#2.
% Q-bar in the typewriter face, like the latch's other signal names.
\newcommand{\Qbar}{\ensuremath{\overline{\mbox{\texttt{Q}}}}}

\newcommand{\hint}[8]{%
  \draw[line width=.9pt,-{Stealth[length=2.8pt,width=2.6pt]}]
    (#1,#2) .. controls (#5,#6) and (#7,#8) .. (#3,#4);}

% The scanned arrows (img/), placed by the matrix they had in the old PDF.
\newcommand{\scan}[7]{%
  \begin{scope}[shift={(current page.south west)},x=1pt,y=1pt,
                cm={#2,#3,#4,#5,(#6,#7)}]
    \node[anchor=south west,inner sep=0,transform shape]
      at (0,0) {\includegraphics[width=1pt,height=1pt]{#1}};
  \end{scope}}

% Cards: a rounded tile in the deck's paper grey.
\newcommand{\card}[4]{\draw[fill=cardbg,draw=cardline,line width=.5pt,
  rounded corners=3pt] (#1,#2) rectangle (#3,#4);}

% A numbered step: a yellow disc with a bold number, centred at #1,#2.
\newcommand{\stepdisc}[3]{%
  \node[circle,fill=hl,minimum size=10.96pt,font=\scriptsize\bfseries,
        text height=5.6pt,text depth=0pt] at (#1,#2) {#3};}

\begin{document}

% ════════════════════════════════════════════════════════════════════════════
% 1  Title
\begin{frame}
\begin{page}
  \ctxt[font=\LARGE]{181.42}{90.68}{Beamer$^+$}
  \ctxt[font=\large]{181.42}{110.29}{Interactive Presentations}
  \ctxt[font=\normalsize,text=cap]{181.42}{146.00}{A Guided Tour}
  \ctxt[font=\footnotesize]{181.42}{169.62}{press \arrowkey{} to begin}
\end{page}
\end{frame}

% ════════════════════════════════════════════════════════════════════════════
% 2  What's inside: twelve tiles, one per feature
\begin{frame}
\begin{page}
  \head{35.20}{what’s inside}
    {Beamer$^+$ turns a static PDF into a live, interactive lecture}
    {every tile below has a demo later in this deck, so press \arrowkey{} to explore}

  \foreach \cx in {10.87,94.95,179.04,263.12}
    \foreach \cy in {85.68,137.48,189.29}
      \card{\cx}{\cy}{\cx+78.70}{\cy+45.83};

  \begin{scope}[font=\scriptsize]
    \ctxt{50.22}{123.33}{video \& audio}   \ctxt{134.31}{123.33}{3D models}
    \ctxt{218.40}{123.33}{split view}      \ctxt{302.48}{123.33}{reveal steps}
    \ctxt{50.22}{175.14}{ink \& shapes}    \ctxt{134.31}{175.14}{magnifier}
    \ctxt{218.40}{175.14}{live widgets}    \ctxt{302.48}{175.14}{Python}
    \ctxt{50.22}{226.95}{live polls}       \ctxt{134.31}{226.95}{scratch paper}
    \ctxt{218.40}{226.95}{screen recording}\ctxt{302.48}{226.95}{notes PDF}
  \end{scope}

  \begin{scope}[line width=.9pt,draw=ink]
    % video & audio
    \draw[rounded corners=1.5pt] (39.26,96.64) rectangle (61.18,110.59);
    \fill[hl] (47.23,99.63) -- (55.20,103.62) -- (47.23,107.60) -- cycle;
    % 3D models: a box, lid on
    \fill[hl] (128.33,100.62) -- (134.31,96.64) -- (140.28,100.62) -- (134.31,104.61) -- cycle;
    \draw[line join=round] (128.33,110.59) -- (128.33,100.62) -- (134.31,96.64) -- (140.28,100.62)
      -- (140.28,110.59) -- (134.31,106.60) -- cycle;
    \draw (128.33,100.62) -- (134.31,104.61) -- (140.28,100.62);
    \draw (134.31,104.61) -- (134.31,106.60);
    % split view
    \draw[rounded corners=1pt] (206.44,96.64) rectangle (217.40,110.59);
    \draw[fill=hl,rounded corners=1pt] (219.39,96.64) rectangle (230.35,110.59);
    % reveal steps
    \draw (291.52,97.63) -- (310.45,97.63);
    \draw (291.52,101.62) -- (306.46,101.62);
    \draw[black!35,dash pattern=on .9pt off 1pt] (291.52,105.60) -- (312.44,105.60);
    \draw[black!35,dash pattern=on .9pt off 1pt] (291.52,109.59) -- (308.46,109.59);
    \fill[hl] (311.44,102.61) -- (315.43,105.60) -- (311.44,108.59) -- cycle;
    % ink & shapes
    \fill[hl,rounded corners=1pt] (38.27,152.43) rectangle (60.18,157.41);
    \draw[rec] (38.27,160.40) .. controls (43.25,147.45) and (47.23,164.39) .. (51.22,153.43)
      .. controls (54.21,147.45) and (58.19,161.40) .. (62.18,150.44);
    % magnifier
    \draw (132.31,153.43) circle (5.98);
    \draw[line width=1.79pt,line cap=round] (136.80,157.91) -- (142.28,163.39);
    \node[font=\fontsize{6}{6}\selectfont\bfseries] at (132.31,153.43) {a};
    % live widgets: a window with a slider
    \draw[rounded corners=1.5pt] (206.44,148.44) rectangle (230.35,162.39);
    \fill[ink] (209.42,150.44) circle (.7) (211.42,150.44) circle (.7);
    \draw (206.44,152.43) -- (230.35,152.43);
    \draw (210.42,157.41) -- (226.36,157.41);
    \draw[line width=.6pt,fill=hl] (220.39,157.41) circle (2.19);
    % Python prompt
    \node[anchor=base west,font=\ttfamily\small] at (292.75,158.40) {>{}>{}>};
    \draw[line width=.4pt,draw=hl!80!black] (309.07,157.98) -- (312.20,157.98);
    % live polls
    \fill[black!45] (38.86,208.22) rectangle (43.65,214.20)
                    (50.82,210.21) rectangle (55.60,214.20)
                    (56.80,205.23) rectangle (61.58,214.20);
    \fill[hl] (44.84,202.24) rectangle (49.62,214.20);
    \draw (38.27,214.20) -- (62.18,214.20);
    % scratch paper: graph paper and a parabola
    \begin{scope}[line width=.4pt,draw=black!30]
      \foreach \gx in {126.34,130.32,134.31,138.29,142.28} \draw (\gx,199.35) -- (\gx,215.09);
      \foreach \gy in {203.24,207.22,211.21} \draw (122.45,\gy) -- (146.16,\gy);
    \end{scope}
    \draw (122.35,199.25) rectangle (146.26,215.19);
    \draw[draw={rgb,1:red,.36;green,.36;blue,.84}] (125.34,202.24) .. controls (130.32,216.18) and (137.30,216.18) .. (143.27,202.24);
    % screen recording
    \draw (218.40,207.22) circle (6.97);
    \fill[rec] (218.40,207.22) circle (4.18);
    % notes PDF
    \draw[line join=round] (296.50,215.19) -- (296.50,199.25) -- (305.47,199.25) -- (309.45,203.24) -- (309.45,215.19) -- cycle;
    \draw (305.47,199.25) -- (305.47,203.24) -- (309.45,203.24);
    \draw[black!45,line width=.6pt] (299.49,206.23) -- (306.46,206.23) (299.49,209.22) -- (306.46,209.22);
    \draw[rec] (299.49,212.70) .. controls (301.5,211.0) and (304.0,214.0) .. (306.46,212.21);
  \end{scope}
\end{page}
\end{frame}

% ════════════════════════════════════════════════════════════════════════════
% 3  Videos / audio — a falling-objects diagram; the video sits on the right
\begin{frame}
\begin{page}
  \head{26.09}{videos/audio}
    {supported formats include \texttt{.mp4}, \texttt{.mov}, \texttt{.mpeg}, \texttt{.mp3}, \texttt{.wav}}
    {example usage: demonstrating Newton’s 2\textsuperscript{nd} law in a Physics course}

  % the vacuum chamber, drawn faint: the video plays over the right half
  \draw[black!12,line width=.8pt,rounded corners=5pt] (27.53,67.80) rectangle (134.79,199.45);
  \draw[black!22,line width=.9pt] (31.18,194.58) -- (131.14,194.58);
  \foreach \hx in {33.62,43.37,...,121.39} \draw[black!14,line width=.5pt] (\hx,194.58) -- (\hx-5.36,198.72);
  \draw[black!18,line width=.6pt,dash pattern=on 2.99pt off 2.99pt] (34.84,104.37) -- (127.48,104.37);
  % the ball and the feather
  \fill[black!55] (52.39,104.37) circle (10.97);
  \fill[white] (48.49,101.20) circle (2.19);
  \draw[black!38,line width=.65pt] (108.95,101.93) -- (108.95,111.69);
  \foreach \fy/\fd in {103.88/2.68, 106.81/2.44, 109.49/2.19}
    \draw[black!22,line width=.5pt] (108.95-4.63*\fd/2.68,\fy-\fd) -- (108.95,\fy) -- (108.95+4.63*\fd/2.68,\fy-\fd);
  % equal arrows
  \draw[black!62,line width=1.89pt,-{Triangle[length=4.3pt,width=6.8pt]}] (52.39,116.07) -- (52.39,147.04);
  \draw[black!62,line width=1.89pt,-{Triangle[length=4.3pt,width=6.8pt]}] (108.95,112.17) -- (108.95,147.04);
  \draw[black!28,line width=.65pt,{Triangle[length=3.8pt,width=3.8pt]}-{Triangle[length=3.8pt,width=3.8pt]}] (65.31,147.04) -- (96.28,147.04);
  \begin{scope}[text=black!55]
    \txt[font=\scriptsize]{34.79}{87.96}{$m_1 = 7\,\mathrm{kg}$}
    \txt[font=\scriptsize]{96.11}{88.32}{$m_2 \approx 0$}
    \txt[font=\scriptsize]{59.22}{132.58}{$m_1 g$}
    \txt[font=\scriptsize]{87.03}{130.39}{$m_2 g$}
  \end{scope}
  \txt[font=\tiny,text=black!38]{75.11}{112.42}{$t{=}0$}
  \txt[font=\tiny,text=black!38]{53.89}{157.47}{same acceleration, $g$}
  \txt[font=\tiny,text=black!28]{58.58}{186.13}{vacuum chamber}

  \ctxt[font=\footnotesize,text=cap]{181.42}{216.83}{gravitational acceleration is}
  \ctxt[font=\footnotesize,text=cap]{181.42}{227.54}{independent of mass}
  \txt[font=\footnotesize]{294.03}{245.04}{try clicking!}
  \scan{img/arrow-curl.png}{-17.07}{-24.38}{34.484}{-24.147}{271.499}{68.377}
\end{page}
\end{frame}

% ════════════════════════════════════════════════════════════════════════════
% 4  Models — blood through the heart; the 3D heart sits on the right
\begin{frame}
\begin{page}
  \head{20.98}{models}
    {supported formats include \texttt{.glpb}}
    {example usage: explaining blood flow through the heart in an Anatomy course}

  \tikzset{chamber/.style={rounded corners=3pt,line width=.4pt}}
  \draw[chamber,draw=black!18,fill=black!5,rounded corners=4pt] (47.85,76.03) rectangle (162.66,92.47);
  \draw[chamber,draw=deoxy!95,fill={rgb,1:red,.91;green,.91;blue,1}] (36.22,113.03) rectangle (92.91,129.47);
  \draw[chamber,draw=deoxy!95,fill={rgb,1:red,.91;green,.91;blue,1}] (36.22,144.35) rectangle (92.91,160.79);
  \draw[chamber,draw={rgb,1:red,1;green,.65;blue,.65},fill={rgb,1:red,1;green,.92;blue,.92}] (117.60,113.03) rectangle (174.30,129.47);
  \draw[chamber,draw={rgb,1:red,1;green,.65;blue,.65},fill={rgb,1:red,1;green,.92;blue,.92}] (117.60,144.35) rectangle (174.30,160.79);
  \draw[chamber,draw={rgb,1:red,.89;green,.78;blue,.89},fill={rgb,1:red,.97;green,.95;blue,.97},rounded corners=4pt] (47.85,175.74) rectangle (162.66,201.42);
  \draw[black!10,line width=.7pt,dash pattern=on 2.99pt off 2.99pt,rounded corners=7pt] (38.18,112.62) rectangle (172.34,164.41);
  \begin{scope}[line width=1.25pt,-{Triangle[length=4pt,width=5pt]}]
    \draw[deoxy] (96.00,92.67) -- (76.24,110.63);
    \draw[deoxy] (64.56,129.67) -- (64.56,140.89);
    \draw[deoxy] (64.56,160.99) -- (83.33,173.71);
    \draw[oxy] (124.48,175.54) -- (143.25,162.82);
    \draw[oxy] (145.95,144.15) -- (145.95,132.93);
    \draw[oxy] (136.69,112.83) -- (116.93,94.87);
  \end{scope}
  \ctxt[font=\scriptsize,text=black!58]{105.26}{86.27}{body}
  \ctxt[font=\scriptsize,text=black!58]{105.26}{186.28}{lungs}
  \ctxt[font=\tiny,text=black!55]{105.26}{195.70}{$\mathrm{O_2} \rightleftharpoons \mathrm{CO_2}$}
  \begin{scope}[font=\fontsize{7}{8}\selectfont]
    \ctxt[text={rgb,1:red,0;green,0;blue,.7}]{64.56}{123.88}{right atrium}
    \ctxt[text={rgb,1:red,0;green,0;blue,.7}]{64.56}{155.20}{right ventricle}
    \ctxt[text={rgb,1:red,.62;green,0;blue,0}]{145.95}{124.56}{left ventricle}
    \ctxt[text={rgb,1:red,.62;green,0;blue,0}]{145.95}{155.88}{left atrium}
  \end{scope}
  \begin{scope}[font=\tiny,text=black!38]
    \ctxt{50.93}{102.20}{superior/inferior}
    \ctxt{59.12}{108.68}{vena cava}
    \txt{138.37}{104.91}{aorta}
    \txt{19.79}{138.73}{tricuspid valve}
    \txt{152.77}{139.31}{mitral valve}
    \txt{17.62}{170.21}{pulmonary artery}
    \txt{147.98}{170.21}{pulmonary veins}
    \node[anchor=base west,font=\tiny\itshape,text=black!22] at (98.54,118.35) {heart};
  \end{scope}

  \ctxt[font=\footnotesize,text=cap]{181.42}{224.57}{blood flow through the heart}
  \ctxt[font=\footnotesize,text=cap]{181.42}{235.53}{(\textcolor{deoxy!95!black}{deoxygenated} blood / \textcolor{oxy!95!black}{oxygenated} blood)}
  \txt[font=\footnotesize]{294.02}{252.72}{try rotating!}
  \scan{img/arrow-curl.png}{-17.07}{-24.38}{34.484}{-24.147}{271.498}{60.697}
\end{page}
\end{frame}

% ════════════════════════════════════════════════════════════════════════════
% 5  Reveal steps — the four steps are numbered on yellow discs, one column,
%    each centred on its line. Each line is one reveal box in the demo config
%    (demo/config/s4.json).
\begin{frame}
\begin{page}
  \head{35.20}{reveal steps}
    {uncover a slide one piece at a time, with no \texttt{\textbackslash pause} or recompiling}
    {example usage: walking through a proof in a Discrete Math class}

  \txt[font=\normalsize]{10.91}{90.79}{\textbf{Claim.}\enspace $\sqrt{2}$ is irrational.}

  % step y = the line's baseline; the disc sits on the line's maths axis
  \foreach \n/\y in {1/115.88, 2/142.79, 3/167.69, 4/192.59}
    \stepdisc{19.60}{\y-3.70}{\n};
  \txt{30.09}{115.88}{Suppose not: $\sqrt{2} = \dfrac{p}{q}$ with $p, q$ integers in \textit{lowest terms}.}
  \txt{30.09}{142.79}{Squaring, $p^2 = 2q^2$, so $p^2$ is even $\;\Rightarrow\;$ $p$ is even: $p = 2k$.}
  \txt{30.09}{167.69}{Then $4k^2 = 2q^2$, so $q^2 = 2k^2$ $\;\Rightarrow\;$ $q$ is even too.}
  \txt{30.09}{192.59}{So $p$ and $q$ are both even, and share a factor of 2.}

  % flush with the step discs' left edge (19.60 - 10.96/2)
  \node[anchor=base west,fill=hl,font=\bfseries,inner xsep=2pt,
        text height=8.56pt,text depth=1.0pt,inner ysep=1pt] at (14.12,217.98) {Contradiction!};
  \txt{96.31}{217.98}{Lowest terms can’t both be even.\,$\square$}

  \txt[font=\footnotesize]{280.0}{250.64}{press \arrowkey{} to reveal!}
  \hint{271.24}{246.83}{256.37}{235.04}{257.29}{246.83}{259.28}{240.86}
\end{page}
\end{frame}

% ════════════════════════════════════════════════════════════════════════════
% 6  Ink & annotations — a projectile problem (the demo's saved ink is drawn
%    over this exact layout)
\begin{frame}
\begin{page}
  \head{35.20}{ink \& annotations}
    {pens, highlighters, shapes, \LaTeX{} text boxes, a spotlight and a magnifier}
    {example usage: working a projectile problem live in a Physics class}

  \txt{11.07}{89.07}{A ball is kicked at $v_0 = 20\,\mathrm{m/s}$, at $\theta = 30^\circ$ above level ground.}
  \txt{11.07}{102.62}{How far away does it land?}
  \txt[font=\small]{11.07}{126.08}{The range is}
  \txt{18.13}{151.56}{$R = \dfrac{v_0^2 \sin 2\theta}{g} = \dfrac{(20)^2 \sin 60^\circ}{9.81}$}

  % the kick: ground, trajectory, launch angle, range
  % (raised 25pt, to sit beside the working rather than under it)
  \begin{scope}[shift={(0,-25)}]
  \draw[black!60,line width=.6pt] (196.27,199.25) -- (353.68,199.25);
  \foreach \gx in {200.25,208.22,...,351.69} \draw[black!35,line width=.4pt] (\gx,199.25) -- (\gx-3.98,203.24);
  \draw[black!60,line width=.9pt,dash pattern=on 2.99pt off 2.99pt]
    plot[domain=0:139.47,samples=60] ({204.24+\x},{199.25-80.48*(\x/139.47)*(1-\x/139.47)});
  \draw[line width=.9pt,-{Stealth[length=3pt,width=2.6pt]}] (204.24,199.25) -- (238.86,179.25);
  \draw[line width=.6pt] (222.17,199.25) arc[start angle=0,end angle=-30,radius=17.93];
  \fill[black!75] (204.24,199.25) circle (2.99);
  \fill[black!25] (343.71,199.25) circle (2.99);
  \draw[black!50,line width=.5pt,{Stealth[length=2pt,width=1.5pt]}-{Stealth[length=2pt,width=1.5pt]}] (204.95,211.21) -- (343.01,211.21);
  \node[fill=white,inner xsep=1.5pt,inner ysep=1pt,font=\footnotesize,text=black!50] at (273.97,211.21) {$R$};
  \txt[font=\small]{227.79}{173.06}{$v_0$}
  \txt[font=\footnotesize]{226.49}{196.15}{$\theta$}
  \end{scope}

  % the fine print the magnifier is for
  \txt[font=\fontsize{3}{3.6}\selectfont,text=black!55]{11.07}{245.6}{Psst, you found the fine print! Scroll while the magnifier is on to zoom anywhere from}
  \txt[font=\fontsize{3}{3.6}\selectfont,text=black!55]{11.07}{249.2}{1.5x to 5x. It works on ink, videos and widgets too.}
  \txt[font=\footnotesize]{217.99}{251.70}{try the magnifier on the fine print!}
  \hint{211.17}{248.37}{196.30}{238.57}{197.22}{248.37}{199.21}{244.39}
\end{page}
\end{frame}

% ════════════════════════════════════════════════════════════════════════════
% 7  Split view — Taylor polynomials of sin x
\begin{frame}
\begin{page}
  \head{19.26}{split view}
    {allows two slides to be shown simultaneously}
    {example usage: explaining Taylor series in a Calculus class}

  \txt{10.91}{68.08}{Any smooth function $f$ can be approximated around $x_0 = 0$ as:}
  \txt{127.23}{98.14}{$\displaystyle f(x) \;=\; \sum_{n=0}^{\infty} \frac{f^{(n)}(0)}{n!}\,x^n$}

  \begin{scope}[shift={(175.53,172.35)},x=28.3465pt,y=28.3465pt]
    \foreach \gx in {-2,-1,1,2} \draw[black!8,line width=.5pt] (\gx,-1.55) -- (\gx,1.55);
    \foreach \gy in {-1,1} \draw[black!8,line width=.5pt] (-2.57,\gy) -- (2.53,\gy);
    \draw[black!30,line width=.8pt,-{Triangle[length=3pt,width=4pt]}] (-2.62,0) -- (2.60,0);
    \draw[black!30,line width=.8pt,-{Triangle[length=3pt,width=4pt]}] (0,-1.60) -- (0,1.56);
    \begin{scope}[font=\tiny]
      \foreach \gx in {-2,-1,1,2} \node[text=black!50,anchor=north] at (\gx,-.07) {$\gx$};
      \node[text=black!50,anchor=east] at (-.07,1) {$1$};
      \node[text=black!50,anchor=east] at (-.07,-1) {$-1$};
    \end{scope}
    \node[anchor=base,font=\scriptsize,text=black!40] at (2.75,-.04) {$x$};
    \node[anchor=base,font=\scriptsize,text=black!40] at (-.08,1.97) {$y$};
    \clip (-2.6,-1.6) rectangle (2.6,1.6);
    \draw[black!72,line width=1.10pt,dash pattern=on 2.99pt off 5.98pt,domain=-2.5:2.5,samples=60]
      plot (\x,{\x-\x^3/6+\x^5/120-\x^7/5040});
    \draw[black!55,line width=1.10pt,domain=-2.5:2.5,samples=60] plot (\x,{\x-\x^3/6+\x^5/120});
    \draw[black!38,line width=1.10pt,dash pattern=on 1.1pt off 1pt,domain=-2.2:2.2,samples=60] plot (\x,{\x-\x^3/6});
    \draw[black!22,line width=1.10pt,dash pattern=on 2.99pt off 2.99pt] (-1.55,-1.55) -- (1.55,1.55);
    \draw[black!70,line width=1.59pt,domain=-2.5:2.5,samples=120] plot (\x,{sin(\x r)});
  \end{scope}
  % legend
  \draw[black!70,line width=1pt] (104.66,132.10) -- (117.41,132.10);
  \draw[black!22,line width=1pt,dash pattern=on 2.99pt off 2.99pt] (104.66,138.33) -- (117.41,138.33);
  \draw[black!38,line width=1pt,dash pattern=on 1pt off 1pt] (104.66,144.57) -- (117.41,144.57);
  \draw[black!55,line width=1pt] (104.66,150.81) -- (117.41,150.81);
  \begin{scope}[font=\tiny,text=black!50]
    \txt{120.6}{134.3}{$\sin x$} \txt{120.6}{140.5}{$T_1$}
    \txt{120.6}{146.7}{$T_3$}    \txt{120.6}{153.0}{$T_5$}
  \end{scope}
  \ctxt[font=\footnotesize,text=cap]{181.42}{239.91}{expansion of $f(x) = \sin(x)$ around $x_0 = 0$}

  % "try using split view!" — circled, by hand
  \fill[hl] (36.09,195.45) -- (72.32,191.00) -- (71.12,181.21) -- (34.89,185.66) -- cycle;
  \node[rotate=7,font=\footnotesize,text=black!70] at (52.0,177.0) {try using};
  \node[rotate=7,font=\footnotesize,text=black!70] at (53.0,188.5) {split view!};
  \node[anchor=south west,inner sep=0] at (14.742,236.294) {\includegraphics[width=77.378pt,height=109.458pt]{img/arrow-loop.png}};
\end{page}
\end{frame}

% ════════════════════════════════════════════════════════════════════════════
% 8  Function plotter — the plotter widget fills the right half
\begin{frame}
\begin{page}
  \head{35.20}{function plotter}
    {2D curves and 3D surfaces you can rotate, zoom and edit live}
    {example usage: classifying critical points in a Multivariable Calculus class}

  \txt{11.07}{88.46}{At a critical point, $\nabla f = \mathbf{0}$. The}
  \txt{11.07}{102.01}{\textit{second-derivative test} looks at}
  \txt{48.03}{126.52}{$D = f_{xx} f_{yy} - f_{xy}^2.$}
  \txt[font=\small]{11.07}{158.76}{For $f(x,y) = a\,x^2 - y^2$ at the origin:}
  \txt[font=\small]{11.07}{173.70}{$D = (2a)(-2) - 0^2 = -4a$}
  \node[anchor=base west,fill=hl,font=\small\bfseries,inner xsep=1.6pt,text height=8.4pt,text depth=2.4pt] at (50.45,190.64) {saddle point};
  \txt[font=\small]{11.07}{190.64}{$a > 0$: a}
  \txt[font=\small]{11.07}{204.58}{$a < 0$: a local maximum}

  \txt[font=\footnotesize]{241.59}{251.70}{try dragging the slider for $a$!}
  \hint{234.77}{248.37}{219.91}{238.57}{220.82}{248.37}{222.82}{244.39}
\end{page}
\end{frame}

% ════════════════════════════════════════════════════════════════════════════
% 9  Widgets — an SR latch from two NOR gates, and its truth table
\begin{frame}
\begin{page}
  \head{35.21}{widgets}
    {supported formats include \texttt{.html}}
    {example usage: explaining latches in a Computer Organization class}

  \txt{10.91}{91.11}{A {\setlength{\fboxsep}{2pt}\kern-2pt\colorbox{hl}{\textbf{latch}}\kern-2pt} stores one bit of state using cross-coupled gate feedback.}

  % two NOR gates, cross-coupled
  \begin{scope}[line width=.5pt]
    \foreach \gy in {142.94,188.87} {
      \draw[fill=black!7.5] (81.27,\gy-15.87)
        .. controls (86.43,\gy-12.68) and (89.09,\gy-2.98) .. (87.21,\gy+5.78)
        .. controls (86.21,\gy+10.48) and (84.03,\gy+14.17) .. (81.27,\gy+15.88)
        .. controls (103.09,\gy+15.88) and (103.09,\gy+7.95) .. (108.55,\gy)
        .. controls (103.09,\gy-7.93) and (103.09,\gy-15.87) .. cycle;
      \draw[fill=black!7.5] (110.19,\gy) circle (1.64);
      \draw (74.72,\gy-7.93) -- (86.65,\gy-7.93) (74.72,\gy+7.93) -- (86.65,\gy+7.93);
      \draw (111.82,\gy) -- (156.64,\gy);
    }
    \draw (44.10,135.00) -- (74.71,135.00);
    \draw (44.10,196.81) -- (74.71,196.81);
    \draw (127.56,142.94) -- (127.56,162.13) -- (68.59,162.13) -- (68.59,180.93) -- (74.71,180.93);
    \draw (136.74,188.87) -- (136.74,169.67) -- (74.71,169.67) -- (74.71,150.88);
    \fill (127.56,142.94) circle (.8) (136.74,188.87) circle (.8);
  \end{scope}
  \txt[font=\footnotesize]{34.07}{138.25}{\texttt{S}}
  \txt[font=\footnotesize]{33.18}{200.06}{\texttt{R}}
  \txt[font=\footnotesize]{160.52}{145.32}{\texttt{Q}}
  \txt[font=\footnotesize]{160.52}{191.87}{\Qbar}
  \ctxt[font=\fontsize{7}{8}\selectfont,text=black!50]{100.0}{224.5}{cross-coupled NOR gates}
  \ctxt[font=\fontsize{7}{8}\selectfont,text=black!50]{100.0}{235.0}{inputs are \textit{active-high}}

  % truth table
  \begin{scope}[line width=.4pt]
    \draw (213.43,133.51) -- (332.10,133.51);
    \draw (213.43,148.16) -- (332.10,148.16);
    \draw (213.43,205.55) -- (332.10,205.55);
    \draw (249.57,133.71) -- (249.57,205.35) (287.20,133.71) -- (287.20,205.35);
  \end{scope}
  \begin{scope}[font=\scriptsize]
    \foreach \y/\s/\r/\q/\nq/\st in {144.80/\texttt{S}/\texttt{R}/\texttt{Q}/\Qbar/state,
                                     159.45/0/0/\texttt{Q}/\Qbar/hold,
                                     173.70/1/0/1/0/set,
                                     187.95/0/1/0/1/reset} {
      \ctxt{222.08}{\y}{\s} \ctxt{240.04}{\y}{\r}
      \ctxt{259.59}{\y}{\q} \ctxt{278.60}{\y}{\nq}
      \txt{293.37}{\y}{\st}
    }
    \ctxt{222.08}{202.19}{1} \ctxt{240.04}{202.19}{1} \ctxt{259.59}{202.19}{0} \ctxt{278.60}{202.19}{0}
    \txt[text=red]{293.37}{202.19}{forbidden}
  \end{scope}
  \ctxt[font=\fontsize{7}{8}\selectfont,text=black!50]{272.77}{223.5}{truth table for an SR latch}
  \ctxt[font=\fontsize{7}{8}\selectfont,text=black!50]{272.77}{234.3}{implemented with NOR gates}
\end{page}
\end{frame}

% ════════════════════════════════════════════════════════════════════════════
% 10  Record & share — three cards
\begin{frame}
\begin{page}
  \head{35.20}{record \& share}
    {the lecture doesn’t have to stay in the room}
    {example usage: posting the recording and annotated notes right after class}

  \foreach \cx in {10.86,127.42,243.99} \card{\cx}{90.66}{\cx+106.60}{188.29};
  \begin{scope}[line width=.9pt,draw=ink]
    \draw (64.16,110.34) circle (7.97);
    \fill[rec] (64.16,110.34) circle (4.98);
    \draw[fill=white,line join=round] (173.75,119.30) -- (173.75,101.37) -- (184.71,101.37) -- (188.69,105.35) -- (188.69,119.30) -- cycle;
    \draw (184.71,101.37) -- (184.71,105.35) -- (188.69,105.35);
    \draw[rec] (176.74,114.32) .. controls (179.5,108.5) and (182.5,119.5) .. (185.71,113.32);
    \draw[fill=white,rounded corners=1pt] (289.32,101.37) rectangle (305.26,119.30);
    \fill[black!45] (296.09,103.36) rectangle (298.48,104.96) (296.09,106.35) rectangle (298.48,107.94)
                    (296.09,109.34) rectangle (298.48,110.93) (296.09,112.33) rectangle (298.48,113.92);
    \draw[fill=hl] (294.90,113.82) rectangle (299.68,118.31);
  \end{scope}
  \begin{scope}[font=\footnotesize]
    \foreach \cx/\h/\a/\b/\c/\d in {
        64.16/{record the tab}/{slides, ink, widgets}/{and sound, saved}/{every second so a}/{crash loses nothing},
        180.72/{notes PDF}/{one handout with}/{your ink, every wid-}/{get as it stands,}/{and the poll results},
        297.29/{one ZIP file}/{slides, media, widgets,}/{ink and poll setup}/{together, ready to}/{reopen next term}} {
      \ctxt[font=\footnotesize\bfseries]{\cx}{131.81}{\h}
      \ctxt[text=cap]{\cx}{144.76}{\a} \ctxt[text=cap]{\cx}{155.72}{\b}
      \ctxt[text=cap]{\cx}{166.67}{\c} \ctxt[text=cap]{\cx}{177.63}{\d}
    }
  \end{scope}
  % Record is at the top of the slide navigator, off the page's left edge
  \txt[font=\footnotesize]{27.0}{247.72}{try the record button!}
  \hint{23.0}{245.0}{5.0}{226.0}{13.0}{246.5}{6.5}{237.0}
\end{page}
\end{frame}

% ════════════════════════════════════════════════════════════════════════════
% 11  Make your own — three numbered steps
\begin{frame}
\begin{page}
  \head{35.20}{make your own}
    {start from the PDF you already have}
    {everything you’ve seen in this deck takes three steps}

  \foreach \cx in {10.86,127.42,243.99} \card{\cx}{93.65}{\cx+106.60}{205.23};
  \foreach \cx/\n in {19.58/1,136.14/2,252.70/3} \stepdisc{\cx}{102.37}{\n};
  \begin{scope}[line width=.9pt,draw=ink]
    % a PDF dropping into a dashed window
    \draw[black!45,dash pattern=on 2.99pt off 2.99pt,rounded corners=1.5pt] (47.22,110.34) rectangle (81.10,134.25);
    \draw[fill=white,line join=round] (58.18,129.27) -- (58.18,112.33) -- (67.15,112.33) -- (71.13,116.31) -- (71.13,129.27) -- cycle;
    \draw (67.15,112.33) -- (67.15,116.31) -- (71.13,116.31);
    \draw[-{Stealth[length=2.2pt,width=2.4pt]}] (64.16,103.36) -- (64.16,108.5);
    \node[font=\fontsize{4}{4}\selectfont\bfseries,text=rec] at (64.65,123.5) {PDF};
    % an editor window with a pencil and a plus
    \draw[fill=white,rounded corners=1.5pt] (164.78,109.34) rectangle (196.66,133.25);
    \draw[black!45,line width=.6pt] (168.77,114.32) -- (186.70,114.32) (168.77,118.31) -- (174.75,118.31);
    \draw[draw=black!60,fill=hl,rounded corners=1pt] (178.73,117.31) rectangle (192.68,129.27);
    \draw[line width=1pt] (185.71,120.30) -- (185.71,126.28) (182.72,123.29) -- (188.69,123.29);
    \draw[fill=white] (193.68,106.36) -- (199.66,112.33) -- (201.65,110.34) -- (195.67,104.36) -- cycle;
    % a player with a record dot
    \draw[fill=white,rounded corners=1.5pt] (281.35,109.34) rectangle (313.23,133.25);
    \fill[ink] (293.30,127.27) -- (304.26,121.29) -- (293.30,115.32) -- cycle;
    \draw[fill=white] (312.23,132.26) circle (4.98);
    \fill[rec] (312.23,132.26) circle (2.99);
  \end{scope}
  \begin{scope}[font=\footnotesize]
    \foreach \cx/\h/\hy/\a/\b/\c/\ty in {
        64.16/{load your slides}/150.73/{drop a PDF onto the}/{window, or a Beamer$^+$}/{ZIP you saved earlier}/163.69,
        180.72/{add the extras}/150.73/{open edit mode to place}/{videos, 3D models,}/{widgets and polls}/163.69,
        297.29/{present}/149.79/{draw, poll and record}/{live, then download}/{the ZIP and your notes}/162.74} {
      \ctxt[font=\footnotesize\bfseries]{\cx}{\hy}{\h}
      \ctxt[text=cap]{\cx}{\ty}{\a} \ctxt[text=cap]{\cx}{\ty+10.96}{\b} \ctxt[text=cap]{\cx}{\ty+21.91}{\c}
    }
  \end{scope}
\end{page}
\end{frame}

\end{document}
